% Part: sets-functions-relations
% Chapter: functions
% Section: kinds

\documentclass[../../../include/open-logic-section]{subfiles}

\begin{document}

\olfileid{sfr}{fun}{kin}
\olsection{Jinis-jinis Fungsi}

\begin{explain}
Bakal migunani yen kita ngenalake panggolongan
kanggo sawetara jinis fungsi kang paling
kerep kita temoni.

Kanggo miwiti, kita bisa ngrembug fungsi
kang nduweni sipat yen saben anggota
kodomain iku nilai saka fungsi kasebut.
Fungsi kaya iki diarani !!{surjective},
lan bisa digambarake kaya ing \olref{fig:surjective}.

\begin{figure}
  \olasset{assets/diagrams/surjective.tikz}
  \caption{Fungsi !!^a{surjective} nduweni
    saben !!{element} kodomain minangka nilai.}
  \ollabel{fig:surjective}
\end{figure}
\end{explain}

\begin{defn}[Fungsi !!^{surjective}]
Fungsi $f \colon A \rightarrow B$ iku
\emph{!!{surjective}} yen lan mung yen
$B$ uga jangkauan~$f$, yaiku kanggo
saben $y \in B$ ana paling ora siji
$x \in A$ kang ndadekake~$f(x) = y$,
utawa kanthi lambang:
\[
  (\forall y \in B)(\exists x \in A)f(x) = y.
\]
Kita nyebut fungsi kaya iki
!!a{surjection} saka $A$ menyang $B$.
\end{defn}

\begin{explain}
Yen arep nuduhake yen $f$ iku
!!a{surjection}, kita kudu nuduhake
yen saben objek ing kodomain $f$
iku nilai $f(x)$ kanggo sawijining
masukan $x$.

Gatekna yen saben fungsi \emph{nuwuhake}
!!a{surjection}.
Amarga, kanggo fungsi $f \colon A \to B$,
tetepna $f' \colon A \to \ran{f}$
kanthi $f'(x) = f(x)$.
Amarga $\ran{f}$ \emph{ditetepake}
minangka $\Setabs{f(x) \in B}{x \in A}$,
fungsi $f'$ iki mesthi
!!a{surjection}
\end{explain}

\begin{explain}
Saben fungsi metakake saben masukan
kang bisa ana menyang siji asil.
Nanging uga ana fungsi kang ora
nate metakake masukan beda
menyang asil kang padha.
Fungsi kaya iki diarani !!{injective},
lan bisa digambarake kaya ing \olref{fig:injective}.
\begin{figure}
  \olasset{assets/diagrams/injective.tikz}
  \caption{Fungsi !!^a{injective} ora nate
    metakake rong argumen kang beda
    menyang nilai kang padha.}
  \ollabel{fig:injective}
\end{figure}
\end{explain}

\begin{defn}[Fungsi !!^{injective}]
Fungsi $f \colon A \rightarrow B$ iku
\emph{!!{injective}} yen lan mung yen
kanggo saben $y \in B$ ana paling akeh
siji $x \in A$ kang ndadekake~$f(x) = y$.
Kita nyebut fungsi kaya iki
!!a{injection} saka $A$ menyang~$B$.
\end{defn}

\begin{explain}
Yen arep nuduhake yen $f$ iku
!!a{injection}, kita kudu nuduhake
yen kanggo saben !!{element}s $x$ lan $y$
saka domain $f$, yen $f(x)=f(y)$,
mula $x=y$.
\end{explain}

\begin{ex}
Fungsi konstan $f\colon \Nat \to \Nat$
kang ditemtokake kanthi $f(x) = 1$
iku ora !!{injective} lan ora !!{surjective}.

Fungsi identitas $f\colon \Nat \to \Nat$
kang ditemtokake kanthi $f(x) = x$
iku !!{injective} lan !!{surjective}.

Fungsi panerus $f \colon \Nat \to \Nat$
kang ditemtokake kanthi $f(x) = x+1$
iku !!{injective} nanging ora !!{surjective}.

Fungsi $f \colon \Nat \to \Nat$
kang ditemtokake kanthi:
\[
  f(x) =
  \begin{cases}
    \frac{x}{2} & \text{yen $x$ genep} \\
    \frac{x+1}{2} & \text{yen $x$ ganjil.}
  \end{cases}
\]
iku !!{surjective}, nanging ora !!{injective}.
\end{ex}

\begin{explain}
Asring kita arep ngrembug fungsi
kang !!{injective} lan !!{surjective}
bebarengan. Fungsi kaya iki diarani
!!{bijective}. Wujude kaya fungsi
ing \olref{fig:bijective}.
!!^{bijection}s kadhangkala uga diarani
\emph{korespondensi siji-menyang-siji},
amarga masangake anggota kodomain
karo anggota domain kanthi tunggal.
\begin{figure}
  \olasset{assets/diagrams/bijective.tikz}
  \caption{Fungsi !!^a{bijective} masangake
    anggota kodomain karo anggota
    domain kanthi tunggal.}
  \ollabel{fig:bijective}
\end{figure}
\end{explain}

\begin{defn}[!!^{bijection}]
Fungsi $f \colon A \to B$ iku
\emph{!!{bijective}} yen lan mung yen
!!{surjective} lan !!{injective}.
Kita nyebut fungsi kaya iki
!!a{bijection} saka $A$ menyang~$B$
(utawa antarane $A$ lan~$B$).
\end{defn}

\end{document}
